\documentclass[preview]{standalone}

\usepackage{amsmath}
\usepackage{amssymb}
\usepackage{stellar}
\usepackage{definitions}

\begin{document}

\id{modules-over-principal-ideal-domains}
\genpage

\section{The PID setting}

\begin{snippet}{pid-module-setting}
    Throughout the following pages, \(A\) denotes a
    \principalidealdomain. In particular, \(A\) is a
    \uniquefactorizationdomain, and its
    \irreducibleelement[irreducible elements] are precisely its
    \primeelement[prime elements].
\end{snippet}

\begin{snippetproposition}{pid-nonzero-prime-ideals-maximal}{Nonzero prime ideals in a PID}
    Every nonzero \primeideal of a \principalidealdomain is a
    \maximalideal. Consequently, if \(I\idealin A\) is nonzero, then
    \[
        I\text{ is prime}
        \ifandonlyif
        A/I\text{ is a \field}.
    \]
    In particular, for every \primeelement \(p\in A\), the quotient
    \(A/(p)\) is a \field.
\end{snippetproposition}

\begin{snippetproof}{pid-nonzero-prime-ideals-maximal-proof}{pid-nonzero-prime-ideals-maximal}{Nonzero prime ideals in a PID}
    Let \(P=(p)\neq(0)\) be prime, and suppose
    \[
        (p)\subseteq I\subseteq A.
    \]
    Since \(A\) is a PID, write \(I=(a)\). Then \(p=ab\) for some
    \(b\in A\). The element \(p\) is prime, hence irreducible, so either
    \(a\) or \(b\) is a unit. In the first case \(I=A\); in the second,
    \(a\) and \(p\) are associates and \(I=P\). Thus \(P\) is maximal.
    The quotient characterization is the standard correspondence between
    prime or maximal ideals and integral-domain or field quotients.
\end{snippetproof}

\end{document}
